\section{Odtwarzanie wyników}
\label{app:odtwarzanie}

Wszystkie wyniki i rysunki można odtworzyć z plików konfiguracyjnych zapisanych w katalogach eksperymentów. Potrzebne są Rust w wersji 1.88.0, instalowany przez \texttt{rustup} na podstawie pliku \texttt{rust-toolchain.toml}, oraz narzędzie uv, które instaluje Pythona 3.13 i zależności w wersjach zapisanych w pliku \texttt{uv.lock}. Zbiory wzorców, także zbiory grupowe i bitmapy zdjęć, są dołączone do programu jako pliki, więc nie trzeba ich tworzyć od nowa.

Całość odtwarza jeden skrypt uruchamiany w głównym katalogu repozytorium (Listing~\ref{lst:polecenia}). Wykonuje on kolejno polecenia zapisane w komentarzach na początku plików konfiguracyjnych wszystkich eksperymentów: najpierw obliczenia w katalogu \texttt{src/compute}, potem rysunki w katalogu \texttt{src/analysis}. Wyniki i metadane (pełna konfiguracja, wersja programu) trafiają do katalogu \texttt{results} obok pliku konfiguracji, a rysunki raportu do katalogu \texttt{figures}. Pomiar czasu i pamięci zależy od komputera, dlatego skrypt wykonuje go tylko z opcją \texttt{-{}-benchmark} (\texttt{-Benchmark} w PowerShellu).

\begin{code}[H]
\begin{Verbatim}[fontsize=\small, commandchars=\\\{\}]
\txtGray{# Linux, macOS lub Git Bash w systemie Windows}
\txtYellow{scripts/reproduce.sh}

\txtGray{# Windows, PowerShell}
\txtYellow{pwsh} -File scripts/reproduce.ps1
\end{Verbatim}
    \caption{Polecenia odtwarzające wszystkie eksperymenty i rysunki.}
    \label{lst:polecenia}
\end{code}

\noindent Jako przykład na Listingu~\ref{lst:konfiguracja} pokazano konfigurację eksperymentu \texttt{cats-noise}. Zawiera ona ziarno, z którego wyprowadza się ziarna wszystkich prób, oba zbiory kotów, obie reguły uczenia, oba tryby dynamiki, poziomy szumu i liczbę prób na wzorzec. Opcja \texttt{stability\_fields} dodatkowo zapisuje pola wyrównane wszystkich wzorców.

\begin{code}[H]
    \lstinputlisting[style=config, firstline=9]{../data/experiments/cats-noise/config.toml}
    \caption{Konfiguracja eksperymentu \texttt{cats-noise} bez komentarzy.}
    \label{lst:konfiguracja}
\end{code}
